\section{Related Work}
\label{sec:related}

\subsection{Frozen-Feature Reference Matching in Few-Shot AD}
\label{sec:related:matching}

소수의 정상 참조를 활용하는 이상 탐지 방법은 패치 매칭, 영상 정합, 특징 복원 및 정상 부분공간 모델링으로 나눌 수 있다.
패치 매칭 계열의 PatchCore~\cite{patchcore}는 사전학습 CNN의 다층 특징을 연결한 정상 메모리 뱅크와 최근접 거리를 사용하며, AnomalyDINO~\cite{anomalydino}와 SuperAD~\cite{superad}는 동결 DINOv2 특징을 활용한다.
RegAD~\cite{regad}와 FastRecon~\cite{fastrecon}는 각각 영상 정합과 특징 복원으로 정상 참조와의 차이를 측정하고, SubspaceAD~\cite{subspacead}는 정상 부분공간에 대한 잔차를 사용한다.
MVTec AD 2의 few-shot 연구~\cite{fewshotad2baseline}도 동결 특징의 계층 선택과 참조 수에 따른 결과를 분석한다.
이러한 정상 모델링 방식에 더해, ScopeAD는 최근접 매칭에서 계층·해상도별 관측을 확장하고 정상 참조 척도와 계층별 가중치를 함께 사용하는 점수 결합에 초점을 둔다.

\subsection{Multi-Scale Observation and Score Fusion}
\label{sec:related:fusion}

다층·다해상도 관측은 서로 다른 수준의 공간 정보를 결합하는 접근이다.
DPT~\cite{ranftl2021vision}는 여러 ViT 계층의 토큰을 학습된 디코더로 결합하고, S$^2$~\cite{shi2024scaling}는 동결 백본의 다중 척도 특징을 타일 처리와 재조립으로 구성한다.
이상 탐지에서는 PatchCore의 특징 연결과 SuperAD의 계층별 이상 맵 평균처럼 특징 또는 점수 단계에서 관측을 결합한다~\cite{patchcore,superad}.
FreePatch~\cite{freepatch}는 정상 참조 정보를 활용하여 패치 매칭 점수를 보정하는 접근이다.
이러한 접근은 관측을 특징 단계에서 합칠지 개별 매칭 후 점수 단계에서 합칠지, 그리고 정상 참조 정보를 결합에 어떻게 사용할지에 따라 구분할 수 있다.
ScopeAD는 계층·입력 해상도별 매칭 거리를 유지하고, 관측별 정상 참조 척도와 계층별 상대 가중치를 함께 사용하는 비대칭 결합으로 이를 통합한다.
구체적인 결합 계수와 정상 참조 추정은 \cref{sec:method:fusion}에서 정의한다.

\subsection{Distribution Shift and Low-FPR Localization}
\label{sec:related:shift}

획득 조건 변화에 대응하는 연구는 정상 외관의 변화와 결함을 구별하는 문제를 다룬다.
AeBAD 연구~\cite{aebad}는 실제 도메인 변화를 포함한 데이터셋과 masked multi-scale reconstruction 방법을 제안하며, MVTec AD 2~\cite{mvtecad2}는 조명 변화와 작은 결함을 포함한 평가 환경을 제공한다.
평가 측면에서 AUPIMO~\cite{aupimo}는 낮은 오경보 허용치에서의 국소화 품질을 측정하여 전체 곡선 지표를 보완한다.
ScopeAD는 정상 서포트로 정한 점수 척도와 고정 결합 규칙을 쿼리에 적용하며, 추가 학습이나 쿼리 적응 없이 분포 변화 조건의 저-FPR 국소화를 평가한다.
이때 정상 거리의 척도 보정과 고정 배포 임계값의 오경보율 보장은 구분한다.
